     5th INTERNATIONAL MATHEMATICS COMPETITION FOR UNIVERSITY
                                  STUDENTS
                 July 29 - August 3, 1998, Blagoevgrad, Bulgaria

                                                  First day

                                      PROBLEMS AND SOLUTIONS


Problem 1. (20 points) Let V be a 10-dimensional real vector space and U1 and U2 two linear subspaces
such that U1 ⊆ U2 , dimIR U1 = 3 and dimIR U2 = 6. Let E be the set of all linear maps T : V −→ V which
have U1 and U2 as invariant subspaces (i.e., T (U1 ) ⊆ U1 and T (U2 ) ⊆ U2 ). Calculate the dimension of E
as a real vector space.
Solution First choose a basis {v1 , v2 , v3 } of U1 . It is possible to extend this basis with vectors v4 ,v5 and
v6 to get a basis of U2 . In the same way we can extend a basis of U2 with vectors v7 , . . . , v10 to get as
basis of V .
    Let T ∈ E be an endomorphism which has U1 and U2 as invariant subspaces. Then its matrix, relative
to the basis {v1 , . . . , v10 } is of the form
                                                                           
                                            ∗ ∗ ∗ ∗ ∗ ∗ ∗ ∗ ∗ ∗
                                         
                                              ∗ ∗ ∗ ∗ ∗ ∗ ∗ ∗ ∗           
                                         
                                              ∗ ∗ ∗ ∗ ∗ ∗ ∗ ∗ ∗           
                                          0 0 0 ∗ ∗ ∗ ∗ ∗ ∗ ∗ 
                                                                           
                                          0 0 0 ∗ ∗ ∗ ∗ ∗ ∗ ∗ 
                                          0 0 0 ∗ ∗ ∗ ∗ ∗ ∗ ∗  .
                                                                           
                                                                           
                                          0 0 0 0 0 0 ∗ ∗ ∗ ∗ 
                                                                           
                                          0 0 0 0 0 0 ∗ ∗ ∗ ∗ 
                                                                           
                                          0 0 0 0 0 0 ∗ ∗ ∗ ∗ 
                                            0 0 0 0 0 0 ∗ ∗ ∗ ∗

So dimIR E = 9 + 18 + 40 = 67.

Problem 2. Prove that the following proposition holds for n = 3 (5 points) and n = 5 (7 points), and
does not hold for n = 4 (8 points).
    “For any permutation π1 of {1, 2, . . . , n} different from the identity there is a permutation π2 such
that any permutation π can be obtained from π1 and π2 using only compositions (for example, π =
π1 ◦ π1 ◦ π2 ◦ π1 ).”
Solution
    Let Sn be the group of permutations of {1, 2, . . . , n}.
    1) When n = 3 the proposition is obvious: if x = (12) we choose y = (123); if x = (123) we choose
y = (12).
    2) n = 4. Let x = (12)(34). Assume that there exists y ∈ Sn , such that S4 = hx, yi. Denote by K
the invariant subgroup
                                 K = {id, (12)(34), (13)(24), (14)(23)}.
    By the fact that x and y generate the whole group S4 , it follows that the factor group S4 /K contains
only powers of ȳ = yK, i.e., S4 /K is cyclic. It is easy to see that this factor-group is not comutative
(something more this group is not isomorphic to S3 ).
    3) n = 5
    a) If x = (12), then for y we can take y = (12345).
    b) If x = (123), we set y = (124)(35). Then y 3 xy 3 = (125) and y 4 = (124). Therefore (123), (124), (125) ∈
hx, yi- the subgroup generated by x and y. From the fact that (123), (124), (125) generate the alternating
subgroup A5 , it follows that A5 ⊂ hx, yi. Moreover y is an odd permutation, hence hx, yi = S5 .
    c) If x = (123)(45), then as in b) we see that for y we can take the element (124).
    d) If x = (1234), we set y = (12345). Then (yx)3 = (24) ∈ hx, yi, x2 (24) = (13) ∈ hx, yi and
 2
y = (13524) ∈ hx, yi. By the fact (13) ∈ hx, yi and (13524) ∈ hx, yi, it follows that hx, yi = S 5 .

                                                       1
   e) If x = (12)(34), then for y we can take y = (1354). Then y 2 x = (125), y 3 x = (124)(53) and by c)
S5 = hx, yi.
   f) If x = (12345), then it is clear that for y we can take the element y = (12).

Problem 3. Let f (x) = 2x(1 − x), x ∈ IR. Define
                                                                     n
                                                          z }| {
                                                     fn = f ◦ . . . ◦f .
                              R1
a) (10 points) Find limn→∞ 0 fn (x)dx.
                          R1
b) (10 points) Compute 0 fn (x)dx for n = 1, 2, . . ..
Solution. a) Fix x = x0 ∈ (0, 1). If we denote xn = fn (x0 ), n = 1, 2, . . . it is easy to see that
x1 ∈ (0, 1/2], x1 ≤ f (x1 ) ≤ 1/2 and xn ≤ f (xn ) ≤ 1/2 (by induction). Then (xn )n is a bounded non-
decreasing sequence and, since xn+1 = 2xn (1 − xn ), the limit l = limn→∞ xn satisfies l = 2l(1 − l), which
implies l = 1/2. Now the monotone convergence theorem implies that
                                                       Z 1
                                               lim           fn (x)dx = 1/2.
                                              n→∞       0


      b) We prove by induction that

                                                                    2n
                                                   1    2n −1      1
(1)                                        fn (x) = − 2         x−
                                                   2               2

holds for n = 1, 2, . . . . For n = 1 this is true, since f (x) = 2x(1 − x) = 12 − 2(x − 12 )2 . If (1) holds for
some n = k, then we have
                                                                     k                                 2k
                           fk+1 (x)   = fk (f (x)) = 21 − 22 −1     1         1 2
                                                                    2 − 2(x − 2 )                − 12
                                                                    k
                                               k                   2
                                      = 21 − 22 −1 −2(x − 12 )2
                                               k+1            k+1
                                      = 21 − 22 −1 (x − 21 )2
which is (2) for n = k + 1.
   Using (1) we can compute the integral,
                     Z 1               "           n                    2n +1 #1
                                   1    22 −1                      1                       1     1
                        fn (x)dx =   x− n                       x−                     =     −         .
                      0            2   2 +1                        2                       2 2(2n + 1)
                                                                                 x=0


Problem 4. (20 points) The function f : IR → IR is twice differentiable and satisfies f (0) = 2, f 0 (0) = −2
and f (1) = 1. Prove that there exists a real number ξ ∈ (0, 1) for which

                                              f (ξ) · f 0 (ξ) + f 00 (ξ) = 0.

Solution. Define the function
                                                            1 2
                                               g(x) =         f (x) + f 0 (x).
                                                            2
Because g(0) = 0 and
                                            f (x) · f 0 (x) + f 00 (x) = g 0 (x),
it is enough to prove that there exists a real number 0 < η ≤ 1 for which g(η) = 0.
     a) If f is never zero, let
                                                   x      1
                                             h(x) = −        .
                                                   2 f (x)


                                                                 2
Because h(0) = h(1) = − 21 , there exists a real number 0 < η < 1 for which h0 (η) = 0. But g = f 2 · h0 ,
and we are done.
    b) If f has at least one zero, let z1 be the first one and z2 be the last one. (The set of the zeros is
closed.) By the conditions, 0 < z1 ≤ z2 < 1.
    The function f is positive on the intervals [0, z1 ) and (z2 , 1]; this implies that f 0 (z1 ) ≤ 0 and f 0 (z2 ) ≥ 0.
    Then g(z1 ) = f 0 (z1 ) ≤ 0 and g(z2 ) = f 0 (z2 ) ≥ 0, and there exists a real number η ∈ [z1 , z2 ] for which
g(η) = 0.
                                         2
Remark. For the function f (x) = x+1        the conditions hold and f · f 0 + f 00 is constantly 0.

Problem 5. Let P be an algebraic polynomial of degree n having only real zeros and real coefficients.
a) (15 points) Prove that for every real x the following inequality holds:
(2)                                             (n − 1)(P 0 (x))2 ≥ nP (x)P 00 (x).

b) (5 points) Examine the cases of equality.
Solution. Observe that both sides of (2) are identically equal to zero if n = 1. Suppose that n > 1. Let
x1 , . . . , xn be the zeros of P . Clearly (2) is true when x = xi , i ∈ {1, . . . , n}, and equality is possible
only if P 0 (xi ) = 0, i.e., if xi is a multiple zero of P . Now suppose that x is not a zero of P . Using the
identities
                                        n
                             P 0 (x) X 1            P 00 (x)    X            2
                                     =            ,          =                            ,
                             P (x)     i=1
                                           x − xi    P (x)           (x − xi )(x − xj )
                                                                             1≤i<j≤n
we find                                  2                   n
                                P 0 (x)             P 00 (x) X n − 1                      X               2
                  (n − 1)                      −n           =                −                                       .
                                P (x)               P (x)     i=1
                                                                  (x − xi )2                      (x − xi )(x − xj )
                                                                                        1≤i<j≤n
But this last expression is simply
                                                                                       2
                                                     X              1        1
                                                                         −                   ,
                                                                  x − xi   x − xj
                                                1≤i<j≤n

and therefore is positive. The inequality is proved. In order that (2) holds with equality sign for every real
x it is necessary that x1 = x2 = . . . = xn . A direct verification shows that indeed, if P (x) = c(x − x1 )n ,
then (2) becomes an identity.

Problem 6. Let f : [0, 1] → IR be a continuous function with the property that for any x and y in the
interval,
                                         xf (y) + yf (x) ≤ 1.
a) (15 points) Show that
                                                          Z 1
                                                                                  π
                                                                      f (x)dx ≤     .
                                                              0                   4
b) (5 points) Find a function, satisfying the condition, for which there is equality.
Solution Observe that the integral is equal to
                                            Z π2
                                                 f (sin θ) cos θdθ
                                                          0
and to                                                  Z π2
                                                                  f (cos θ) sin θdθ
                                                          0
So, twice the integral is at most
                                                              Z π2
                                                                               π
                                                                       1dθ =     .
                                                                  0            2
                  √
Now let f (x) =       1 − x2 . If x = sin θ and y = sin φ then
                            xf (y) + yf (x) = sin θ cos φ + sin φ cos θ = sin(θ + φ) ≤ 1.


                                                                         3
